
égalités 2.4.2 sont des polynômes en @@M0@@, à coefficients rationnels, et il suffit donc, pour montrer leur égalité, de montrer qu'ils prennent même valeur pour toutes valeurs entières de @@M1@@. Mais les relations obtenues par substitution de valeurs entières à @@M2@@ sont celles qui expriment que le pré-@@M3@@-anneau @@M4@@ (2.5.a)) est un @@M5@@-anneau; elles sont donc satisfaites.

\textbf{Définition 2.7.} \uline{On appelle} anneau binômial \uline{un anneau vérifiant les conditions} 2.6.i) \uline{et} ii), \uline{muni de la structure de} @@M6@@-\uline{anneau définie par} 2.6.1.

\textbf{Exemple 2.8.} Soient @@M7@@ un site, @@M8@@ un anneau binômial, @@M9@@ le faisceau associé au préfaisceau constant de valeur @@M10@@ sur @@M11@@ ; alors @@M12@@ est un anneau binômial.

En effet, les conditions 2.6 se conservent par passage à la limite inductive.

En particulier, @@M13@@ est un anneau binômial.

2.9. Soient @@M14@@ un anneau binômial, @@M15@@ une @@M16@@-algèbre. Pour toute série @@M17@@, et tout @@M18@@, on peut définir une série @@M19@@ de la façon usuelle : on écrit @@M20@@, et on pose :
@@M21@@
la dernière égalité ne pouvant être écrite que modulo torsion. Ces formules sont obtenues par substitution à partir des formules analogues pour @@M22@@ ; les relations suivantes, propriétés classiques de l'exponentielle, restent donc encore vraies pour @@M23@@, @@M24@@ :
@@M25@@



% Exposé V continuation macros
\providecommand{\Ghat}{\widehat G}
\providecommand{\Gg}{\widehat G_g}
\providecommand{\Gcirc}{\widehat G_{\circ}}
\providecommand{\Gnabla}{\widehat G_{\nabla}}
\providecommand{\Grhat}{\widehat{\operatorname{Gr}}}
\providecommand{\Log}{\operatorname{Log}}
\providecommand{\Fil}{\operatorname{Fil}}
\providecommand{\eps}{\varepsilon}
\renewcommand{\id}{\operatorname{id}}
\providecommand{\Sym}{\operatorname{Sym}}
\providecommand{\Mon}{\operatorname{Mon}}


